\chapter{犯罪形态}

% Source PDF physical page 116; see audit/source_coverage.json.
\begin{EditionOpeningPanel}
1.犯罪形态是指犯罪在时间上呈现的形态。一般教材只讲犯罪未完成形态（犯罪预备、未遂和中止），不讲犯罪既遂。但由于犯罪既遂是法考的五星级考点，因此本讲将犯罪既遂一并讲解，因此本讲标题是犯罪形态。

2.过失犯罪只有成立与否问题，无犯罪形态，无既遂未遂问题。

3.犯罪形态是犯罪的一种终局性形态，而非暂时性停顿。即使是犯罪预备、犯罪未遂、犯罪中止等犯罪未完成形态，也是一种终局性形态。因此，各犯罪形态之间是相互排斥关系，不能并存。即在同一犯罪中，只能出现一个犯罪形态。

\InsertDiagram{figures/ch11_states.tex}

4.形态分布图：

\InsertDiagram{figures/ch11_timeline.tex}
\end{EditionOpeningPanel}

\section{犯罪预备}

[问题]狗蛋欲杀小芳，想买一把刀，没钱买刀，便在超市打工挣钱。打工挣钱，是否属于犯罪的预备行为？（打工买刀案）

第22条为了犯罪，准备工具、制造条件的，是犯罪预备。

对于预备犯，可以比照既遂犯从轻、减轻处罚或者免除处罚。

犯罪预备是一种犯罪的终局性形态。犯罪预备的成立条件有三个：第一，主观上为了着手实行犯罪；第二，客观上实施了预备行为；第三，由于意志以外原因未能着手实行犯罪。

% Source PDF physical page 117; see audit/source_coverage.json.
\subsection*{一、主观上为了着手实行犯罪}

第22条第1款中的“为了犯罪”应当理解为“为了着手实行犯罪”。因此，为了预备犯罪而做的准备，不是犯罪预备行为。例如，为了实行抢劫而购买凶器的行为，是预备行为；为了购买凶器打出租车前往五金商店的行为，不是犯罪的预备行为。同理，本节开头“打工买刀案”中，打工行为不属于犯罪的预备行为。

\subsection*{二、客观上实施了预备行为}

预备行为，是指对法益已经造成一定危险的行为，危险尚未达到紧迫的程度。

[注意1]\CJKunderwave{犯罪预备行为与犯意表示的区分}：关键看对法益的危险程度。预备行为对法益造成一定危险：犯意表示对法益无任何危险，只是将犯意单纯流露于外部。例如，小王对室友说：“我讨厌小张，我想放他血！”这只是犯意流露。如果小王进而买了一把刀准备实施，就是犯罪的预备行为。

[注意2］预备犯、犯罪预备和预备行为的区分。预备犯是指构成犯罪预备的罪犯；犯罪预备就是在预备阶段因意志以外原因未能着手实行的犯罪形态；预备行为是指实施了为实行犯罪而做的准备行为。犯罪预备与预备行为是两个不同概念。

\subsection*{三、由于意志以外原因未能着手实行犯罪}

预备阶段有两个犯罪形态：犯罪预备和犯罪中止。二者的区分：关键在于自动性。犯罪预备是由于意志以外的原因而被迫放弃犯罪，而犯罪中止是自动放弃犯罪。

例如（2010年第5题），妻子乙欲毒杀丈夫甲，做了有毒饭菜，等待甲回家，期间，乙决定不等了，打消杀人念头，将饭菜倒掉。乙构成预备阶段的犯罪中止。

\section{犯罪未遂}

[问题］某日，狗剩驾驶推土机去往工地，看到前方一位女士（小美）向前走，顿起奸淫之心。狗剩将推土机的斗轻轻放下来，从身后将小美铲进斗里，驾驶推土机一路狂奔，行至僻静处，放下斗，准备奸淫。小美爬出来，狗剩一看，又不想强奸了，便连忙道歉。小美骂骂咧咧离去。对狗剩如何处理？（推土机强奸案）

第23条已经\CJKunderwave{着手实行}犯罪，由于犯罪分子\CJKunderwave{意志以外的原因}而\CJKunderwave{未得逞}的，是犯罪未遂。

对于未遂犯，可以比照既遂犯从轻或者减轻处罚。

犯罪未遂的成立条件有两个：第一，已经着手实行犯罪；第二，因为意志以外原因未得逞。“未得逞”，表明不是犯罪既遂。“因为意志以外原因”，表明不是犯罪中止。

\subsection*{一、犯罪未遂与犯罪预备}

\subsubsection*{（一）区分标准}

\CJKunderwave{二者的区分}：犯罪预备处在预备阶段，犯罪未遂处在实行阶段。预备阶段与实行阶段

% Source PDF physical page 118; see audit/source_coverage.json.
的分界点是着手。因此，关键区分在于着手，着手之前是犯罪预备，着手之后是犯罪未遂。预备阶段的犯罪中止与实行阶段的犯罪中止，区分标准也是是否着手。

\CJKunderwave{着手的判断标准}：行为对法益是否造成现实、紧迫、直接的危险。若是，则属于着手。

例如，甲欲毒杀乙，向乙邮寄毒药，寄出时不是着手，收到时是着手。\footnote{这种主流标准也称为\CJKunderwave{实质的}客观结果（状态）说。理论上还有三种少数说。第一种观点（主观说）：犯意表现出来时就是着手。例1，为了入户盗窃而打破他人家窗户玻璃，就是盗窃罪的着手。例2，为了入户抢劫，携带凶器进入他人家院子，就是着手。第二种观点（形式的客观说）：开始实施刑法分则规定的构成要件行为就是着手。例如，为了实施保险诈骗，制造保险事故，就是保险诈骗罪的着手。第三种观点（实质的客观行为说）：开始实施具有现实危险性的行为时，就是着手。例如，甲为了杀乙，给乙寄出毒药时，就是着手。}（2019年试题）

\subsubsection*{（二）易考情形}

1.抢劫罪：对人使用暴力胁迫或其他强制手段时是着手。例1（2006年第54题），甲乙合谋抢劫出租车，准备好凶器，上了出租车，谎称去郊外，打算在郊外动手。出租车行驶到某检查站，甲乙发现检查人员怀疑自己，便逃离。甲乙是犯罪预备。例2（2024年试题），甲欲抢劫哨兵手中的枪，夜晚埋伏于哨兵身后草丛，没有机会下手，遂放弃，于返回途中被抓捕。甲构成抢劫枪支罪的犯罪预备。

2.强奸罪：对妇女使用暴力胁迫或其他强制手段时是着手，而非奸淫时。例1，将妇女向车里拖，准备在车里强奸，向车里拖时就是着手。例2，在本节开头“推土机强奸案”中，狗剩用推土机铲的行为，属于已经着手。例3，甲欲强奸乙女，埋伏在乙女下班回家路上，由于乙提前路过，甲没等到。甲构成犯罪预备。（2006年第54题）

3.盗窃罪。入户盗窃，翻院墙，不算着手；翻进去后，若屋里没人，撬门算着手；若屋里有人，进门后才算着手。如果题干没有特别交代，就以屋里没人为准。

4.诈骗罪：为了诈骗而伪造证件的行为是预备行为，开始向被害人实施诈骗才是着手。打电话诈骗，话术铺垫后，开始要钱时是着手（多数说）。（2024年主观题）

5.保险诈骗罪：为了保险诈骗而制造事故的行为是预备行为，开始向保险公司提出索赔才是着手。制造完事故后，向保险公司打电话询问索赔事宜，不是着手。

6.诬告陷害罪：为了诬告陷害而写诬告材料是预备行为，向有关机关告发才是着手。

\subsubsection*{（三）特殊问题}

1.隔离犯的着手。隔离犯是指行为与结果之间存在时空上的间隔。例如，邮寄毒药杀人。隔离犯的着手标准：看邮寄物在途中有无危险，如果有，则寄出时为着手；如果没有，则收到打开时为着手。

例1，甲邮寄一盒有毒饼干给乙，乙收到时，甲才算杀人着手。如果收到但不打开，那么打开时算着手。（2019年试题）

例2，甲邮寄炭疽热病毒粉给乙，由于该病毒粉在邮寄途中一旦泄漏就有传播危险，所以寄出时就是杀人着手。

2.间接正犯的着手：以被利用人为标准，被利用人着手，就视为间接正犯着手。例如，甲指使小孩入室盗窃乙家，甲指使时不算着手，小孩着手实施时才算着手；小孩在去乙家路上被人拐卖了，则甲构成盗窃罪间接正犯的犯罪预备。

% Source PDF physical page 119; see audit/source_coverage.json.
\subsection*{二、未遂犯与不能犯}

不能犯，是指行为人虽然主观有犯意，但是客观行为不具有任何法益侵害危险，所以无罪。例如（稻草人案），甲在荒无人烟的沙漠里误将稻草人当作仇人乙而开枪。甲不构成犯罪。理由是，犯罪行为必须是个危害行为，危害行为必须是对法益制造危险的行为。甲的开枪行为对任何人的生命法益都没有制造危险，不是危害行为，因此不构成犯罪。

不能犯的法律效果是作无罪处理，因此不能犯也常被称为不可罚的不能犯。\footnote{提示：不能犯作无罪处理，是就不能犯的行为举止而言的。如果行为人还有其他预备行为，且对法益产生了一定危险，则能构成犯罪预备。例如，甲欲杀仇人乙，拿着枪四下寻找乙，来到沙漠，将稻草人当作乙而开枪。甲向稻草人的开枪行为属于不能犯，无罪。但甲拿着枪寻找乙，这种寻找行为对乙的法益产生了一定危险，属于预备行为，构成犯罪预备。如果题中没有特别交代，则不需要考虑犯罪预备的情形。}

\subsubsection*{（一） 分类}

1.对象不能犯：因为不存在犯罪对象，导致不可能构成犯罪。例如，在沙漠里将稻草人当作仇人射击。

[注意］对象错误与对象不能犯。二者的相同点是行为人主观上都存在对象认识错误。区别在于：对象错误中，行为具有危险性，属于危害行为，构成犯罪；对象不能犯中，行为不具有危险性，不属于危害行为，也即行为不构成犯罪。

例1，甲看到前方有个人，误以为是仇人李四，实际是王五，向其开枪，没有击中。这种行为对王五有危险，甲构成故意杀人罪未遂。这属于对象错误，一般在“事实认识错误”中讲解。

例2，甲在四下无人的沙漠里，误将稻草人当作仇人而开枪。这种行为对他人生命没有任何危险，甲无罪。这属于对象不能犯，一般在“犯罪未遂”中讲解。

2.手段不能犯：因为手段不可能产生任何危险，导致不可能构成犯罪。例如，用一把无法射击的坏枪射杀人。

[注意1]打击错误与手段不能犯。二者的相同点是行为客观上都出现了方法错误。区别在于：打击错误中，手段具有危险性，属于危害行为，构成犯罪；手段不能犯中，手段不具有危险性，不属于危害行为，不构成犯罪。

例1，甲举枪射杀乙，因为没有瞄准，打死了乙身边的丙。甲的行为具有危险性，构成犯罪。这属于打击错误，一般在“事实认识错误”中讲解。

例2，甲举枪射杀乙，发现枪原来是一把早已生锈的坏枪，根本无法射击。甲的行为不具有危险性，不构成犯罪。这属于手段不能犯，一般在“犯罪未遂”中讲解。

[注意2］迷信犯。迷信犯是手段不能犯的一种，也无罪。这是指由于迷信无知，误以为手段会产生结果，其实手段没有任何危险性，根本不会产生结果。例如，捏小人插针诅咒对方死亡。又如，误以为寺庙的香灰能毒死人而用来投毒杀人。

% Source PDF physical page 120; see audit/source_coverage.json.
\subsubsection*{（二）不能犯与未遂犯的区分}

\InsertDiagram{figures/manual_table_p120_1.tex}

二者的区分标准：行为对法益有无危险。如果有，就是未遂犯；如果没有，就是不能犯。

\InsertDiagram{figures/manual_p120_2.tex}

1.主观说与客观说

（1）主观说（少数说）：行为有无危险，以行为人的主观认识为准。行为人主观上认为有危险，那么行为就有危险。

（2）客观说（多数说）：行为有无危险，是一种客观判断，以客观现实情况为准。

例如，甲将面粉交给乙，谎称是毒品，让乙运输，乙以为是毒品而运输，途中被抓。主观说认为，从乙主观认识看，运输行为有侵害法益的危险性，乙构成运输毒品罪未遂。客观说认为，从客观现实情况看，甲欺骗了乙，乙根本不可能运到毒品，运输行为没有侵害法益的危险性。乙不构成运输毒品罪，属于对象不能犯。

2.客观危险说与具体危险说

客观说包括客观危险说和具体危险说。

（1）客观危险说（少数说）：行为有无危险，应从事后角度和科学专家角度判断。这是“事后诸葛亮”的标准。

（2）具体危险说（多数说）：行为有无危险，应从行为时一般人的角度判断。

例如，甲欲杀乙，向乙投放1毫克毒药，乙服用后未死。事后，专家鉴定查明，该毒药1毫克不会致人死亡，没有致命危险。客观危险说认为，基于此，甲的行为没有危险，不构成故意杀人罪，属于手段不能犯。具体危险说认为，从行为时一般人的角度看，1毫克毒药也是毒药，万一甲在行为时投放量大一点，就有致命危险。因此，甲构成故意杀人罪，属于犯罪未遂。

具体危险说的判断方法：从行为时一般人角度看问题，就是用辩证的眼光看问题，也即用联系的、发展的、全面的眼光看问题。从事后诸葛亮角度看问题，就是用孤立的、静止的、片面的眼光看问题（形而上学的眼光看问题）。如果用联系的、发展的、全面的眼光看，行为对法益有危险，但由于偶然因素没有导致实害结果，此时不能因为没有实害结果就否认行为的危险性。危险性就是指可能性。判断危险性，就是判断可能性。偶然没发生，不代表没有发生的可能性。

% Source PDF physical page 121; see audit/source_coverage.json.
例1，甲拦路抢劫乙，发现乙身无分文。不能因为乙身无分文而认为抢劫行为没有危险性。甲构成抢劫罪未遂。

例2，甲夜晚潜伏乙家谋杀乙，从窗户外向乙的床铺猛开枪，实际上床上无人。乙当时在洗手间。不能因为床上无人就认为开枪行为没有危险性。朝一个有人居住的家里开枪，对居住人制造了紧迫危险。甲构成故意杀人罪未遂。

例3，甲向乙开枪射击，乙因为身穿防弹衣而毫发未损。不能因为乙没有受伤而认为开枪行为没有危险性。甲构成故意杀人罪未遂。（2009年第52题）

[巩固练习]题1，甲乙相距100米，甲向乙开枪，没有射中乙。事后鉴定发现，甲的手枪最大射程只有90米，不会射中乙。不能因为这一次没有打中就认为这种行为没有危险性。甲构成故意杀人罪未遂。

题2，甲欲杀乙，向躺在床上的乙开枪。情形一，实际上，开枪前，乙已经死了有二十分钟。甲的行为没有侵害乙的生命的危险性，属于对象不能犯，不构成故意杀人罪。情形二，实际上，开枪前两分钟，乙刚死。从联系、发展的眼光看，人咽气有个过程，有可能早死几分钟或晚死几分钟，而甲对此又不知道，甲有可能早几分钟开枪。因此，甲的行为具有侵害乙的生命的危险性、可能性，所以构成故意杀人罪未遂。（2003年第4题）

题3，甲入室盗窃，发现女主人睡着了，便强奸，发现是男人。由于有妇女存在的可能性，因此甲有侵害妇女的危险性，构成强奸未遂。假如题中特别交代“该住处不可能有妇女存在，是个和尚宿舍”，则属于对象不能犯，不构成强奸罪。（2005年第7题）

题4（欺骗案），甲欺骗乙，故意给乙一把空枪，让乙开枪打死丙。乙不知情而开枪。乙的行为对丙没有危险性，不构成故意杀人罪。甲也不构成故意杀人罪的教唆犯。

[总结】对行为危险性的判断公式：从客观角度\ensuremath{\rightarrow}行为时一般人的角度（联系、发展、全面的眼光）\ensuremath{\rightarrow}行为有无危险。\hypertarget{extended-note-121-1-ref}{}\textsuperscript{\hyperlink{extended-note-121-1}{附注}}

\section{犯罪中止}

[问题］狗剩夜晚要拦路强奸过路陌生妇女，将妇女二丫打晕后，认为二丫颜值太低当场离去。对狗剩如何处理？（拦路强奸案）

第24条在犯罪过程中，\CJKunderwave{自动放弃}犯罪或者\CJKunderwave{自动有效地防止犯}罪结果发生的，是犯罪中止。

对于中止犯，\CJKunderwave{没有造成损害的}，应当免除处罚；\CJKunderwave{造成损害的}，应当减轻处罚。

\par\noindent\hypertarget{extended-note-121-1}{\hyperlink{extended-note-121-1-ref}{\textbf{附注}}}\quad 提示：概念界定问题。传统理论将不能犯分为“可罚的不能犯”与“不可罚的不能犯”。但是，张明楷教授认为，不能犯就是指不可罚的不能犯。传统理论将未遂犯分为“能犯未遂”与“不能犯未遂”。张明楷教授认为这种分法没有必要。参见张明楷：《刑法学》（第七版），第522、532页。实际上，这是概念界定问题。对应的事实特征和结论都是一样的。在此不存在实质的观点之争，只是概念界定之争。概念界定之争没有实际意义。

\InsertDiagram{figures/manual_table_p121_1.tex}

% Source PDF physical page 122; see audit/source_coverage.json.
\subsection*{一、成立条件一：自动性}

\subsubsection*{（一）犯罪中止与犯罪未遂的区分}

区分标准（弗兰克公式）：“能达目的而不欲”，是中止；“欲达目的而不能”，是未遂。能达目的而不欲，是\CJKunderwave{指能够继续犯罪的前提下}，\CJKunderwave{自动放弃犯罪}。欲达目的而不能，是指\CJKunderwave{无法继续犯罪的前提下}，\CJKunderwave{被迫放弃犯罪}。

“能够继续犯罪”，属于前提条件（外在条件）。“自动放弃犯罪”，属于主观条件（内在条件）。做题时，需先判断“继续犯罪的可能性”，再判断“放弃犯罪的自动性与被迫性”。抛开前者，孤立地判断后者没有意义，也无法判断。例如，男生甲苦追女生乙未果，对室友讲：“我主动放弃。”孤立地看，甲是自动放弃，是中止，但前提是甲追不到，所以属于被迫放弃，是未遂。

1.前提条件（外在条件）：具有继续犯罪的可能性

关于“能不能继续犯罪”的判断，理论上曾有纯粹主观说和纯粹客观说。这两说均不妥当。理论上现采取客观说。

（1）纯粹主观说认为，应根据犯罪人主观看法来判断能不能继续犯罪，犯罪人认为不能继续犯罪，就认定不能继续犯罪。然而，该说并不妥当。例如，在本节开头“拦路强奸案”中，依主观说，狗剩认为不能继续犯罪，那么便认定为不能继续犯罪。依此，狗剩便是在不能继续犯罪的情况下被迫放弃犯罪，构成未遂。这种结论并不妥当。

（2）纯粹客观说认为，应根据客观的、物理的条件来判断能不能继续犯罪，从客观角度看能继续犯罪，就认定能继续犯罪。依照该说，在本节开头“拦路强奸案”中，在客观上能继续强奸，在此前提下放弃犯罪，属于主动放弃，应认定为中止。该说在此似乎没问题。但是该说仍存在缺陷。例如，甲欲强奸陌生妇女，忽然发现对方是自己的亲妹妹，便放弃。依纯粹客观说，在客观上对亲妹妹仍能够继续强奸，在此前提下放弃便属于主动放弃，应认定为中止。然而，这种结论并不妥当。对此应认定为未遂。

（3）客观说（多数说）认为，依据\CJKunderwave{社会一般人的看法}，看在当时情境下能不能继续犯罪。社会一般人的标准（平均值、中间值）是个客观标准，因此称为客观说。

例1，本节开头“拦路强奸案”中，根据社会一般人的看法，可以继续强奸。此时放弃，成立犯罪中止。

例2，甲本想拦路强奸过路的陌生妇女，打晕对方后，发现对方是亲妹妹，便放弃。根据社会一般人的看法，无法强奸亲妹妹。此时放弃，成立犯罪未遂。有人认为，他心目中的“社会一般人”能够强奸亲妹妹。然而，他心目中的“社会一般人”已经不是人了，是个连人性都没有的人。“社会一般人”至少是有人性的人。

例3，甲有了婚外情，计划杀死妻子，举刀砍杀妻子，此时孩子走进来。甲不忍心在孩子面前杀妻子，便放弃犯罪。“社会一般人”这个标准应用到具体案件中，也可能存在不同意见。第一种意见认为，从社会一般人的角度看，面对孩子，甲无法继续杀妻子，此时放弃，成立犯罪未遂。第二种意见认为，从社会一般人的角度看，面对孩子，甲仍能够继续杀妻子，此时放弃，成立犯罪中止。第一种意见将“社会一般人”理解成极负责任感的人。然而，“社会一般人”就是大街上的普通人，就是路人甲，道德水准上，既不是极负责任感的人，也不是没有人性的人。就这种人而言，已经决定杀妻子了，此时孩子并不能成为现实障碍。甲只是一时可怜孩子而放弃，属于主动放弃，而非被迫放弃。所以，第

% Source PDF physical page 123; see audit/source_coverage.json.
一种意见属于少数意见，第二种意见属于多数意见。（2019年客观题，2021年主观题）\footnote{提示：有同学认为：“既然例2中，根据社会一般人看法，甲无法强奸亲妹妹（无法侵害近亲属），构成未遂，那么例3中，根据社会一般人看法，甲也应无法杀害妻子（无法侵害近亲属），也应构成未遂”。这种推理忽视了例2与例3的案件事实存在重大区别。例2中，甲的计划目标是陌生妇女，未想到是亲妹妹。例3中，甲的计划目标就是妻子。在此，需要判断的问题是，根据社会一般人看法，在孩子面前，能否杀害妻子。“根据社会一般人看法，无法侵害近亲属”，这个结论的前提是，计划的侵害目标是陌生人，未料遇到近亲属。}

2.主观条件（内在条件）：具有放弃犯罪的自动性，而非被迫性

（1）主观说（多数说）。成立犯罪中止，要求犯罪人在主观上是自动放弃犯罪，而非被迫放弃犯罪。理论上将这种主观心理条件称为主观说（多数说）。如何判断这种主观心理？如果犯罪人是在能够继续犯罪的前提下放弃犯罪，则推定犯罪人是自动放弃犯罪。如果犯罪人是在无法继续犯罪的前提下放弃犯罪，则推定犯罪人是被迫放弃犯罪。例如，甲用刀杀乙，看到乙满脸是血的样子，感到惊愣，便放弃杀乙。依据一般人的看法，甲可以继续犯罪。此时放弃，推定是自动放弃，定中止。

（2）限定主观说（少数说）。该说认为，在主观说的基础上应加个限定条件，也即犯罪人自动放弃犯罪必须是出于真诚悔过的动机，才能成立犯罪中止。例如，甲男嫌妇女长得丑而放弃强奸，由于缺乏悔过动机，不成立犯罪中止。\footnote{关于主观心理的要求，理论上还有犯罪人理性说（少数说）。该说认为，如果犯罪人经过理性冷静权衡后，认为无法继续犯罪，那么表明确实无法继续犯罪，此时放弃，属于被迫放弃，成立未遂。例如，甲不忍心在孩子面前杀害妻子，放弃犯罪。甲放弃犯罪不是经过理性冷静权衡，而是一种感性的心理（同情心、怜心），因此不能认为是无法继续犯罪而放弃犯罪，不能认为是犯罪未遂。该说不足在于，会导致越理性冷酷的犯罪人，越容易成立犯罪中止，

例如，甲是冷酷的黑社会分子，欲强奸过路妇女，发现是亲妹妹，甲经过理性冷酷的权衡，认为能强奸亲妹妹，但是此时不想强奸，放弃强奸。根据犯罪人理性说，甲构成犯罪中止。该结论欠妥当。}该说不足在于，不当缩小了犯罪中止的成立范围。

[总结］关于中止与未遂的常考情形：

1.计划侵害陌生人，未料遇到熟人。结论：\CJKunderwave{遇到普通熟人，放弃，是中止；遇到近}\CJKunderwave{亲属，放弃，是未遂}。

例1，想强奸过路妇女，发现是亲妹妹，放弃，成立未遂；发现是个普通认识的人，放弃，成立中止。

例2，想抢劫过路行人，发现是自已父亲、亲哥哥，放弃，成立未遂；发现是个普通熟人，放弃，成立中止。（2003年第2题）

2.害怕处罚。结论：\CJKunderwave{害怕当场被抓（以当场很可能被抓为前提），此时放弃，是未遂；}\CJKunderwave{害怕日后被抓（不会当场被抓），此时放弃，是中止}。（2004年第4题，2010年第57题）

范例：甲举枪要打死十米外的乙。第一，多名警察出现，举枪对准甲，要求甲放下手枪。甲收枪逃离。第二，乙声称“我爸是公安局局长，一定会抓住你，判你死刑！”甲害怕受到处罚，放弃。第一种情形，害怕当场被抓，此时放弃，是未遂。第二种情形，害怕日后被抓，此时放弃，是中止。

[巩固练习］题1（2024年试题），甲对路边一女士心生列意欲强奸，将其骗到一无人处欲下手，该女士发现远处走来行人，向其大声呼救。甲见状遂放弃。甲害怕当场被抓而被迫放弃，构成犯罪未遂。

题2，甲欲盗窃银行，晚上进入银行后。第一，看到金库门口有防卫激光线，一碰触，警铃就会响，感觉无法突破，强闯会被抓，遂放弃。甲属于未遂。第二，看到金库门口上

% Source PDF physical page 124; see audit/source_coverage.json.
方有摄像头，继续的话会被摄像头拍下来，日后一定会被抓住，遂放弃。甲属于中止。

\subsubsection*{（二）认识错误}

标准案例（脚步声案）：甲入室盗窃，忽然听到门外脚步声，以为主人要回家了，赶紧从阳台逃离。实际上是邻居回家的脚步声。甲成立未遂还是中止？（2006年第54题）

行为人对前提条件“能不能继续犯罪”产生认识错误，由此影响主观条件“认为自已是自动放弃犯罪还是被迫放弃犯罪”。\CJKunderwave{判断标准（主观说）：看主观，主观怎么想，就}\CJKunderwave{怎么定}。例如，上述脚步声案中，甲主观上以为只能被迫放弃犯罪，认为只能未遂。结论便是犯罪未遂。

[巩固练习]题1，甲投毒杀乙，乙呕吐不止，甲又心生怜悯，送乙到医院后其康复。事后鉴定，甲投放的毒药不足以致命，不送医院也死不了。甲主观上以为可以继续而主动放弃犯罪，主观上认为自已在实施犯罪中止。结论便是犯罪中止。（2003年第42题）

题2，甲想杀乙，将乙推下过街天桥，看到乙十分痛苦，又开车送乙到医院抢救。实际上，乙未受致命伤，不送医院也死不了。甲主观上认为自己在实施犯罪中止，就定犯罪中止。（2016年第53题）

[注意］欺骗产生的认识错误。

1.乙欺骗甲，甲产生认识错误。要看甲是哪种性质的认识错误。如果对前提条件“能否继续犯罪”有认识错误，则属于真正重要的认识错误，看主观，主观怎么想，就怎么定。如果对前提条件“能否继续犯罪”没有认识错误，则属于不重要的认识错误，不需要考虑这种认识错误。

例1，甲抓住妇女乙的手腕，欲绑架乙后出卖。乙欺骗甲：“我有性病，不会有人要。”甲信以为真，放弃。甲关于“性病”有认识错误，但甲在前提条件“能够继续实施拐卖妇女罪”上没有认识错误，仍能够继续实施，只是觉得卖不到好价钱而主动放弃，定中止。（2004年第4题）

例2，甲强奸乙女，乙女欺骗甲：“我有性病（或艾滋病）。”甲信以为真，放弃。甲关于“性病、艾滋病”有认识错误，导致甲在前提条件“能否继续犯罪”上有认识错误，认为无法继续犯罪。这种属于真正重要的认识错误。看主观，主观怎么想就怎么定。甲主观认为无法继续，被迫放弃，所以定未遂。

2.乙欺骗甲，甲产生认识错误。但是，甲并没有放弃犯意，犯罪未呈现终局形态。这种认识错误不重要，不需要考虑。

例如，甲强奸乙女，乙女欺骗甲：“我不反抗了，我也认了，但能不能歇会，喘口气？”甲信以为真，便松手，两人坐在路边。2分钟后，乙终于看到行人，大声呼救。甲见状逃离。甲松手，不表示甲主动放弃了犯意，只是暂时性停顿，不能将此举动认定为中止。最后甲被迫逃离，定未遂。

\subsubsection*{（三）特定对象不存在}

这是指行为人的目标是特定对象，该特定对象在客观上不存在，行为人主观上认为只能被迫放弃犯罪；如果特定对象有存在的可能性，或行为对同类对象产生紧迫危险，则定未遂。

% Source PDF physical page 125; see audit/source_coverage.json.
1.特定物不存在

例如，甲盗窃银行的保险柜，好不容易打开后，发现里面只有2元钱，非常失望，转身离去。甲构成盗窃未遂。甲的盗窃对象是特定的财物（巨额现金），因为不存在，只能是未遂。即使他拿走了2元钱，因为相对于巨额现金而言，拿走2元钱跟没拿一样，仍是盗窃未遂。

2.特定人不存在

例1，甲让乙杀丙。乙在大街上发现丙，正要开枪时，突然发现对方不是丙，赶紧收枪。（1）由于特定对象不存在，乙只能构成未遂。（2）定未遂，而非犯罪预备，是因为乙的行为对眼前这个人已经制造了紧迫危险。（3）此时，不能认为：乙客观上可以继续杀人，所以定中止。因为此时乙的行为对象是特定的，特定的行为对象不存在，只能认为是无法继续。（4）当然，如果乙将认识错误进行到底并杀了人。这便是同一构成要件内的对象认识错误，成立杀人既遂。这时定既遂，是因为乙已经造成实害结果了。而本案中的乙在最后一刻没有产生认识错误。

例2，甲抓住杀父仇人乙，举刀要砍杀。情形一（2009年第5题），乙欺骗甲：“我只是执行者，主谋是丙，你应该找他！”甲相信，放掉乙，去找丙。情形二，乙欺骗甲：“不是我杀了你父亲，是丙杀了你父亲！你应该找他！”甲相信，放掉乙，去找丙。情形一中，甲在前提条件“能不能继续杀乙”上没有认识错误，认为能继续杀乙，在此前提下放弃杀乙，属于主动放弃，定中止。情形二中，甲在前提条件“能不能继续犯罪”上有认识错误，认为特定对象不存在，不能继续杀乙，在此前提下放弃杀乙，属于被迫放弃，定未遂。

\subsection*{二、成立条件二：中止行为}

这是指成立犯罪中止，要求在客观上行为人实施了中止行为。根据是否实行终了，犯罪中止分为两种情形。

\subsubsection*{（一）行为实行终了与行为未实行终了}

这里的行为是指实行行为。实行行为实行终了，是指能够导致既遂结果的行为已经实行完毕，否则属于未实行终了。

[提示】行为实行终了不等于犯罪既遂，犯罪既遂要求实害结果发生。

例1，甲举刀砍了乙脖子一刀，导致致命伤。此时属于实行终了。第一，此时甲被抓捕，乙被救活。甲属于实行终了的未遂。第二，此时甲后悔，将乙救活。甲属于实行终了的中止。

例2，甲举刀准备砍乙时，此时属于未实行终了。第一，此时甲被抓捕，甲属于未实行终了的未遂。第二，此时甲后悔，自动放弃。甲属于未实行终了的中止。

\subsubsection*{（二）行为未实行终了的犯罪中止}

行为未实行终了时，只要自动放弃犯罪，就可以成立中止。

1.这里的自动放弃，要求是真实彻底的放弃，而非暂时的停顿。例如，甲入室盗窃，发现财物过多，便出门去叫同伙一起搬运，不是盗窃中止。

2.如果实施财产犯罪，转换犯罪对象不算犯罪中止。例如，甲欲盗窃乙家电冰箱，入室后看到笔记本电脑，便放弃电冰箱而窃得笔记本电脑，是盗窃既遂。

\subsubsection*{（三）行为实行终了的犯罪中止}

行为实行终了后，要成立中止，仅自动放弃犯罪是不行的，还要求有效防止犯罪结果

% Source PDF physical page 126; see audit/source_coverage.json.
发生。“有效防止”包括可能的有效性和实际的有效性。可能的有效性，是指防止措施要有防止结果发生的可能性。实际的有效性，是指防止措施实际防止了结果发生。注意，这种防止措施还要求真诚努力地去完成。

例1，甲砍了乙两刀，导致致命伤，看到乙流血不止又后悔，给乙伤口放了些纸巾便离去，邻居将乙救活。甲的防止措施不具有防止结果发生的可能性，不成立犯罪中止。

例2，甲点燃仓库后又后悔，大喊：“着火了，救火啦！”然后逃离。甲的中止行为没有真诚努力去完成，不成立犯罪中止。

\InsertDiagram{figures/manual_p126_1.tex}

\subsection*{三、成立条件三：有效性}

这里的有效性，是就“实行终了的中止”而言的，是指实际的有效性。即使行为人自动放弃或积极努力防止，但结果仍发生了，也不能成立犯罪中止。

\subsubsection*{（一）没有发生危害结果}

1.没有发生危害结果，是指没有发生行为人追求或放任的、行为性质决定的危害结果，而不是指没有发生任何结果。因此，犯罪中止可以分为没有造成任何危害结果的犯罪中止和造成一定危害结果的犯罪中止。对这两种中止的处罚是不同的（参见下文）。

例如（2003年第42题），王某欲杀死李某，捅了两刀，看到李某流血不止，又心生怜悯，将其送到医院，抢救过来，但仍身受重伤。王某仍成立故意杀人罪的犯罪中止。

2.因果关系问题。犯罪中止的模型如下：犯罪行为一中止行为（防止措施）\ensuremath{\rightarrow}结果未发生。

注意：成立犯罪中止，不要求中止行为（防止措施）与结果未发生之间具有因果关系，也即不要求结果未发生是中止行为的功劳。

例1，甲投毒杀乙，乙疼痛难忍，甲又后悔，立即打120急救电话，急救车马上到来，甲出门迎接急救车，接到急救车，进屋发现乙不在。原来，邻居丙开车将乙送往医院，抢救脱险。甲成立犯罪中止。

例2，甲投毒杀乙，乙疼痛难忍，甲又后悔，送乙去医院。经查明，毒药不可能致人死亡，不用送医也不会死。（1）甲的抢救措施与乙的“未死亡”之间没有因果关系，但即使如此，仍认定甲构成故意杀人罪中止。（2）从认识错误角度也可以分析。本案中甲实际上产生了认识错误，主观上以为能够既遂而主动放弃，实际上客观无法既遂。对这种认识错误，应看主观而定，定中止。（2003年第42题，2016年第53题）

\subsubsection*{（二）发生了危害结果}

一般而言，发生了危害结果，那么行为人就构成犯罪既遂，不构成犯罪中止。但是，如果犯罪行为与危害结果之间没有因果关系，则行为人不构成犯罪既遂，而构成犯罪中止。这种情形被称为有效性的例外。

行为模型：犯罪行为\ensuremath{\rightarrow}中止行为（防止措施）\ensuremath{\rightarrow}介入因素一发生实害结果。

判断犯罪行为与危害结果之间有无因果关系，使用“介入因素两步走”标准来判断。

% Source PDF physical page 127; see audit/source_coverage.json.
第一步，介入因素若异常，则与先前犯罪行为是独立关系。第二步，看谁对结果的发生所起的作用大。如果介入因素阻断了救助行为，则表明介入因素的作用大，那么介入因素与结果有因果关系。（2026年试题）

标准案例（2014年第53题）：甲以杀人故意将乙刺成重伤，又后悔，便开车送乙去医院，在路途中甲不慎酿成车祸，车祸直接导致乙死亡。第一步，介入因素是车祸，比较异常。第二步，车祸阻断了救助行为，并且直接导致死亡结果，对死亡结果的发生作用大，与死亡结果有因果关系。而先前犯罪行为与死亡结果没有因果关系，因此，甲不构成故意杀人罪既遂。只剩下犯罪未遂与中止可供选择。由于未遂是指想继续杀死，因为意志以外原因未能杀死，所以不成立未遂，只能成立犯罪中止。车祸是由甲造成的，甲构成交通肇事罪。最终，对甲定故意杀人罪中止和交通肇事罪，数罪并罚。常见的错误是，对甲只定一个故意杀人罪既遂。

[巩固练习]题1（2015年第6题），甲用毒蛇咬乙，乙中毒，甲后悔，开车送乙去医院，途中乙不想活了，跳车逃走，三小时后中毒死亡。后查明，只要当时送医就不会死。第一步，介入因素是乙的逃离行为，比较异常。第二步，介入因素阻断了救助行为，对死亡结果的发生作用大，与死亡结果有因果关系。而先前犯罪行为与死亡结果没有因果关系，因此，甲不构成故意杀人罪既遂，只能成立犯罪中止。

题2（2020年试题），宋某杀害刘某，致刘某重伤昏迷，生命垂危。宋某心生怜悯，想要抱起刘某送去医院救治，不料脚下一滑，和刘某一起摔倒在地，刘某原本已经生命垂危，加上摔倒，很快死亡。假设一：摔倒行为比较异常。先前重伤行为与摔倒行为对死亡结果的发生，作用都大，二因一果。假设二：摔倒行为不异常，那么先前重伤行为与死亡结果有因果关系。因此，无论摔倒行为是否异常，宋某均构成故意杀人罪既遂。这表明宋某的中止行为没有实际有效性，不能成立犯罪中止。

[总结］中止的有效性\footnote{从“犯罪结果是否发生”也可以总结如下：

1.犯罪结果发生

（1）定既遂。条件：实行行为与结果有因果关系。

（2）定未遂。条件：一是实行行为与结果没有因果关系；二是没有可能有效的中止行为。

（3）定中止。条件：一是实行行为与结果没有因果关系；二是有可能有效的中止行为。

2.犯罪结果未发生

（1）定未遂。条件：没有可能有效的中止行为。

（2）定中止。第一种情形：有可能有效的中止行为，并且，结果未发生是中止行为的功劳。第二种情形：有可能有效的中止行为，但是，结果未发生不是中止行为的功劳。第三种情形，只有自动放弃，没有可能有效的中止行为（防止措施），结果未发生不是自动放弃的功劳。例如，甲乙一起进入丙家盗窃，甲后悔，放弃离开，没有阻止乙。乙也后悔，放弃离开。乙成立中止。甲虽然没有阻止乙，但也成立中止。}

\InsertDiagram{figures/manual_table_p127_1.tex}

% Source PDF physical page 128; see audit/source_coverage.json.
\ContinuedTable{figures/manual_table_p128_1.tex}

\subsection*{四、犯罪中止的处罚}

第24条第2款对于中止犯，没有造成损害的，应当免除处罚；造成损害的，应当减轻处罚。

\subsubsection*{（一）损害结果的认定范围\protect\footnotemark}\footnotetext{参见张明楷：《中止犯中的“造成损害”》，载《中国法学》2013年第5期。}

1.这里的损害结果不是指既遂结果，否则就构成犯罪既遂了。

2.这里的损害结果是刑法要处罚的危害结果。

例1，甲制服了妇女乙欲强奸，猥亵完了后又不想强奸了。由于刑法处罚强制猥亵妇女这种危害结果，所以甲虽成立强奸罪中止，但属于“造成损害的，应当减轻处罚”，而不能认定为“没有造成损害的，应当免除处罚”。

例2，甲入户伤害乙，乙哀求，甲放弃。由于刑法处罚非法侵入住宅的行为，所以甲虽成立故意伤害罪中止，但属于“造成损害的，应当减轻处罚”。（2023年试题）

例3，甲抢劫乙，导致乙轻微伤后，又不想抢劫了。由于刑法不处罚轻微伤这种危害结果，所以甲成立抢劫罪中止，属于“没有造成损害的，应当免除处罚”。

\subsubsection*{（二）损害结果的因果关系}

犯罪中止中有两个行为：一是先前的犯罪行为，二是中止行为（防止措施）。\CJKunderwave{注意}：\CJKunderwave{这}\CJKunderwave{里的“损害结果”是指先前犯罪行为导致的损害结果，而不包括防止措施导致的损害结果}

例1，甲晚上盗窃仓库，将一个零部件从机器上拆下来，导致机器被毁，甲抱着零部件往外走，尚未走出仓库门，又后悔，放下零部件，走出仓库门。甲成立盗窃罪的犯罪中止。甲的先前盗窃行为导致机器被毁，造成了损害结果，应当减轻处罚。\footnote{甲的一个盗窃行为同时触犯盗窃罪和故意毁坏财物罪，想象竞合，择一重罪（重罚）论处。比如，按照盗窃罪，定犯罪中止，减轻处罚，处罚效果如果轻于故意毁坏财物罪，那么最终应以故意毁坏财物罪论处。就故意毁坏财物罪而言，成立犯罪既遂，不是犯罪中止，因为已经将机器毁坏了。}

例2，甲晚上盗窃仓库，拿起地上的一个零部件往外走，尚未走出仓库门，又后悔，便返身打算放回去，不慎脱手，砸死了熟睡的保安。甲成立盗窃罪的犯罪中止。该损害结果（死亡结果）不是先前盗窃行为导致的，而是后面的中止行为导致的，因此，在处罚盗窃罪的犯罪中止时，不属于“造成损害的，应当减轻处罚”，而属于“没有造成损害的，应当免除处罚”。然后再单独处罚过失致人死亡罪。

例3，甲投毒杀乙，药性尚未发作，甲又后悔，便送乙去医院，路上甲不慎酿成车祸，致乙死亡。甲成立故意杀人罪中止。该死亡结果不是先前杀人行为导致的，所以该犯罪中止不属于“造成损害的，应当减轻处罚”，而属于“没有造成损害的，应当免除处罚”。甲酿成车祸，致乙死亡，构成交通肇事罪。对该罪进行单独处罚。

% Source PDF physical page 129; see audit/source_coverage.json.
\section{犯罪既遂}

\subsection*{一、既遂的认定}

\subsubsection*{（一）发生阶段}

既遂结果出现在实行阶段，是由实行行为导致的。如果预备行为偶然导致实害结果，不属于既遂结果，不构成犯罪既遂。

例1，甲欲杀死同事乙，将毒酒放在自己办公桌上，准备晚上给乙喝，然后出门。乙来到甲办公室，不知情竟喝了毒酒，中毒死亡。这属于结果的提前发生。甲的杀人行为只是预备行为，不构成故意杀人罪既遂，而构成故意杀人罪犯罪预备，同时触犯过失致人死亡罪，想象竞合，择一重罪论处。

例2，甲欲抢劫前面独自行走的妇女，便远远地跟踪。妇女发觉有人跟踪，为了“丢卒保车”，将提包扔下便跑。甲捡到提包，还追上妇女并扇一耳光，骂道：“你什么意思？施舍吗？”遂拿着提包离去。甲的抢劫行为只是预备行为，不构成抢劫罪既遂，而构成抢劫罪犯罪预备。虽然妇女放弃了所有权，但是甲属于不当得利，甲有返还义务，不返还，构成侵占罪。对抢劫罪（犯罪预备）和侵占罪既遂并罚。

\subsubsection*{（二）因果关系}

既遂所要求的实害结果必须与实行行为有因果关系，否则不构成犯罪既遂。也即，若因果链条断裂，则不构成犯罪既遂。注意，总论中“因果关系”一讲讲述的因果关系判断方法，主要是针对伤亡结果。对于财产犯罪及其他犯罪的因果关系，应根据该罪的行为结构判断。总结常考情形：

1.诈骗罪。诈骗罪的行为逻辑结构是：欺骗行为一对方陷人认识错误一对方基于认识错误而处分财物一行为人因此取得财物。如果缺少第三步，则即使取得财物，也是未遂。

例1，甲对胡某实施诈骗行为，胡某识破骗局，但胡某觉得甲穷困潦倒，实在可怜，就给其3000元。胡某不是基于认识错误而处分财物。甲成立诈骗罪未遂。（2004年第4题，2013年第61题）

例2，甲欲诈骗保险公司，以汽车被盗为由向保险公司索赔，保险公司报警。警方安排保险公司向甲“理赔”，甲领到保险金后，被埋伏在现场的警察抓获。由于保险公司不是基于认识错误而处分财物，甲构成保险诈骗罪未遂。（2009年第15题）

2.敲诈勒索罪。敲诈勒索罪的行为逻辑结构是：恐吓行为一对方产生恐惧心理一对方基于恐惧心理而处分财物一行为人因此取得财物。如果缺少第三步，则即使取得财物，也是未遂。

例1，甲恐吓乙，如果不给1万元就放乙儿子的血，乙实际上是黑社会老大，不但没有恐惧反而很欣赏甲的勇气，便给了甲1万元。乙不是基于恐惧心理而处分财物。甲成立敲诈勒索罪未遂。

例2，甲敲诈勒索乙的钱财，让乙将钱放到指定地点，否则将乙的丑闻曝光。乙报警后，警察让乙按照甲的指示去放钱，以此诱捕甲。乙将钱放到指定地点，甲刚拿到钱后便被埋伏的警察抓捕。乙不是基于恐惧心理而处分财物。甲是敲诈勒索罪的未遂。

% Source PDF physical page 130; see audit/source_coverage.json.
3.抢劫罪。抢劫罪的行为逻辑结构是：暴力胁迫行为一压制对方反抗一对方无法反抗而交付财物一行为人因此取得财物。如果缺少第三步，则即使取得财物，也是未遂。

例如，甲为抢劫而殴打章某，章某逃跑，甲随后追赶。章某在逃跑时钱包不慎从身上掉下，章某不知情而跑离现场。甲拾得钱包后离开。章某不是因无法反抗而被迫放弃财物。缺少第三步，甲的暴力行为和取得财物之间没有因果关系，甲是抢劫罪未遂。甲的取财行为构成侵占罪。两罪并罚。（2007年第1题）

4.强奸罪。强奸罪的行为逻辑结构是：强制行为一压制妇女反抗一妇女因为无法反抗而被奸淫。

例如，夜晚，狗剩入户强奸小美，小美挣扎中打开灯，发现是同村狗剩哥，便斥问为何这样，狗剩惭愧退出，小美又喊道：“其实我早就喜欢你！”两人遂发生性关系。狗剩成立强奸罪中止，而非既遂。

\subsection*{二、各犯罪形态的联系}

\subsubsection*{（一）终局形态}

犯罪预备、未遂、中止、既遂四个\CJKunderwave{犯罪形态是终局性的结束，不是暂时性的停顿}。就同一起犯罪而言，如果出现了一种犯罪形态后，就不可能再出现其他犯罪形态。例如，出现未遂后，就不可能再出现中止或既遂。

[结论］\CJKunderwave{就同一起犯罪而言，不可能并存两个犯罪形态，犯罪形态之间是排斥关系，}\CJKunderwave{而非并存关系。终局性形态的成立条件：}第一，客观上，犯罪行为彻底结束；第二，主观上，犯意彻底消除。

\subsubsection*{（二）同一起犯罪中的排斥关系（五星级考点）}

1.既遂排斥中止

犯罪既遂后，不要把事后悔过行为或返还行为当成中止行为。例如，甲盗窃到同学的一台笔记本电脑，放进自己抽屉，又后悔，便悄悄送还，仍是既遂。

2.既遂排斥未遂

犯罪既遂后，财物又被其他人拿走的，不再成立未遂。例如，甲入户盗窃到乙的箱子，扔到院墙外僻静处，翻墙出来找时没找到，原来被路过的丙拿走了。甲将箱子扔到僻静处时就已经既遂，不能因为后来没找到而变成未遂。

3.中止排斥未遂

（1）第一个时间节点呈现终局性形态。

例如，甲早上入室强奸妇女乙，乙为了摆脱，谎称：“我其实也喜欢你，但此刻身体不舒服，你晚上来，好吗？”甲信以为真，便离去。晚上甲提着为乙买的药品来到乙家，被埋伏的警察抓获。甲早上离开现场后愿意等待大半天晚上再来，表明甲心里真以为乙喜欢自己，甲的强奸的犯意便消除了，心里带着通奸的想法。因此，甲早上离开现场时，犯罪呈现终局性形态，构成犯罪中止，晚上被捕就不能再认定为犯罪未遂。

（2）第一个时间节点未呈现终局性形态。

例如（2016年第53题），甲持刀抢劫乙，用刀威胁乙。乙声称：“把刀放下，我给钱就是了！”甲便准备收刀。乙趁其不备，用脚踢掉刀，然后将甲制服。甲收刀时，由于犯罪行为没有结束，还要继续要钱，犯意也没有消除，故未呈现终局性形态，此时不构成犯罪中止。被制服时呈现终局性形态，构成未遂。

% Source PDF physical page 131; see audit/source_coverage.json.
4.未遂排斥中止

（1）第一个时间节点呈现终局性形态。

例如，甲欲杀乙，猛砍10刀，乙身受重伤，甲认为乙很快必死无疑，扔掉刀便离开现场。2小时后，甲为了寻找、藏匿刀回到现场，发现乙还没死，又可怜乙，将其送到医院救活。对此需先判断终局性时刻是哪一刻？客观上，甲离开现场，犯罪行为彻底结束。主观上，在甲离开现场几分钟后，甲认为自已的杀人计划已经实现，此时杀人的犯意便消除了。此时犯罪便呈现终局性形态，由于实际上人没死，所以是犯罪未遂。一旦成立未遂，事后的中止行为便不能成立犯罪中止。\footnote{引申的可能情形一：上述关于未遂与中止的分析，存在一个前提，也即前面的杀人行为没有导致被害人死亡。只有人最终没死，才会启动上述分析。假如乙被送往医院，因伤势过重不治身亡。由于重伤到死亡有个发展过程，最终死亡，且死亡结果与杀人行为有因果关系，则应认定为杀人既遂。此时，便不需要启动上述关于未遂与中止的分析。

引申的可能情形二：上述真题案例中，假如甲2小时后为了找刀回到现场，发现乙还没死，便继续砍，砍死乙。甲先是杀人未遂，回来后又砍，属于新的犯罪，构成杀人既遂。由于这是针对同一个对象的同种犯罪行为，属于连续犯，只定一个故意杀人罪既遂即可。}（2003年第42题，2022年试题）

（2）第一个时间节点未呈现终局性形态。

例1，甲砍乙10刀，乙重伤昏迷。甲停止并观望。10分钟后，乙苏醒过来，请求甲不要杀自己。甲将乙救活。在砍完10刀时，犯罪没有出现终局性形态。甲最终成立犯罪中止。

例2，丙为了杀害丈夫赵某，在赵某睡觉时将煤气阀门打开，然后出门在院子里等待。半小时后，丙估计时间差不多了，回家查验，发现赵某还没死，还在抽搐，顿生悔意，将赵某救活。丙出门等待时，由于犯意没有消除，所以犯罪并没有出现终局性形态。丙最终成立犯罪中止。\footnote{提醒：许多考生只注意犯罪形态之间的区分，例如，预备与未遂、未遂与中止等，而漏掉“犯罪形态之间的排斥关系”这个考点。实际上，这个考点很重要，而且考查方式更隐蔽，应当注意。}

\subsubsection*{（三）重复侵害行为}

这是指连续重复实施相同的侵害行为。问题是，第一个行为是否呈现终局形态？判断标准：第一个行为与第二个行为在时间上是否具有紧密联系性。一般而言，如果相隔一天或一夜，则认为第一个行为终局了。

例如，甲向乙开枪，有五发子弹，第一枪没打中，准备打第二枪，见乙可怜，放弃。由于第一枪和第二枪之间时间间隔很短，所以第一枪结束时犯罪没有呈现终局形态，因此，对第一枪没打中不能认定为犯罪未遂。此时应当整体考察，第二枪呈现终局形态，定犯罪中止。

［巩固练习］题1，甲第一天晚上盗窃，没能撬开保险柜，放弃。次日晚上，拿着合适工具又来撬，中途自动放弃。第一个行为与第二个行为隔了一夜，表明第一个行为呈现终局形态，构成未遂。第二个行为构成中止。不过不需要并罚，哪个处罚更重定哪个，也即按照未遂论处。

题2，甲在轿车内欲杀害妻子乙，用啤酒瓶砸头，由于空间受限，砸了几次，无法砸死，便放下酒瓶，用手掐乙脖子，正在掐时，看乙可怜，放弃。放弃酒瓶未呈现终局形态，不定未遂。放弃掐脖子，定中止。

% Source PDF physical page 132; see audit/source_coverage.json.
\OriginalPanel{典型真题}

下列哪个行为构成犯罪中止？（2022年回忆版金题）\footnote{[答案]A项，该案件属于结果的提前发生。高某投毒四次，表明其杀人已经着手，因此死亡结果属于既遂结果。高某构成故意杀人罪既遂，而非犯罪中止。A项不入选。B项，甲捅了数刀，离开现场，此时甲的犯罪行为已经结束，犯意已经消除，表明犯罪已经呈现终局形态，具体而言是犯罪未遂，因为乙实际上未死。甲的抢救行为不能成立犯罪中止，只能视为事后悔过行为。B项不入选。C项，制造毒品罪的既遂标准是制造出来毒品。李某构成制造毒品罪既遂，而非犯罪中止。C项不入选。D项，宋某的敲诈行为已经终了，要成立犯罪中止，需要采取防止措施，也即防止王某遭受财产损失。宋某采取了防止措施，也即向王某打电话告知。虽然王某仍然打了钱，但宋某及时退还，防止了王某遭受实际的财产损失。因此，宋某构成犯罪中止。D项入选。本题答案为D。}

A.高某欲杀害王某，计划投毒六次以达致死量，致使王某中毒死亡，投毒四次后，由于害怕而放弃继续投毒。但前四次投毒便已导致王某死亡

B.丈夫甲杀害妻子乙，捅了数刀，以为乙死亡，便离开现场。后甲回到现场，发现乙未死，可怜乙，将乙送往医院抢救成功

C.李某欲制造毒品，在化学实验室制造出80g毒品，后因害怕便将毒品扔进水里

D.宋某敲诈勒索王某，要求王某给自己汇款30万元。后宋某害怕坐牢，便告知王某：“我放弃了，不用打钱了，请放心。”但是，王某基于害怕，仍旧给宋某账户汇入30万元。宋某得知后及时退还

